\documentclass{article}

\usepackage{amsmath}
\usepackage{amsthm}
\usepackage{amssymb}
\usepackage{hyperref}
\usepackage{listings}
\usepackage{xcolor}

\lstset{
    language=Python,
    basicstyle=\ttfamily\footnotesize,
    keywordstyle=\color{blue},
    commentstyle=\color{green!60!black},
    stringstyle=\color{red},
    showstringspaces=false,
    breaklines=true,
    frame=single,
    numbers=left,
    numberstyle=\tiny\color{gray},
    tabsize=4,
    columns=flexible,
    keepspaces=true,
    basewidth=0.5em
}

\usepackage{enumitem}
\setlist[enumerate,1]{label=\arabic*.}
\setlist[enumerate,2]{label=(\alph*)}
\setlist[enumerate,3]{label=\roman*.}

\title{Written HW4}
\author{
    Logan Tanner \\ 
    I use Neovim (BTW) \\ 
    \\
    Full HW repository at \\
    \href{https://github.com/LoganCTanner/School/tree/main/Math/425-Probability}{github.com/LoganCTanner/School}
}
\date{\today}

\begin{document}
\maketitle

\clearpage
\begin{enumerate}

    \item   Suppose that $ X \sim normal(\mu,\sigma^2) $. Find the following probabilities:
        \begin{enumerate}

            \item   $P\{|X - \mu | \leq \sigma\}$
                \\Answer:
                \begin{align*}
                    P\{|X - \mu | \leq \sigma\} &= P\{-\sigma \leq X - \mu \leq \sigma\} \\
                    &= P\{-1 \leq \frac{X - \mu}{\sigma} \leq 1\} \\
                    &= P\{\frac{X - \mu}{\sigma} \leq 1\} - \left( 1 - P\{\frac{X - \mu}{\sigma} \leq 1\}\right) \\
                    &= 2P\{\frac{X - \mu}{\sigma} \leq 1\} - 1 \\
                    &= 1.6826 - 1 = \fbox{0.6826}
                \end{align*}
                Given that $\frac{X-\mu}{\sigma} = Z \sim \text{normal}(0,1)$ and the values in the table provided.
            \item   $P\{|X - \mu | \leq 2\sigma\}$ \\
                Answer:\\
                It follows from part (a) that \\
                $ P\{|X - \mu | \leq 2\sigma\} = 2P\{\frac{X - \mu}{\sigma} \leq 2\} - 1 = \fbox{0.9544} $
            \item   $P\{|X - \mu | \leq 3\sigma\}$\\
                Answer:\\
                $ P\{|X - \mu | \leq 3\sigma\} = 2P\{\frac{X - \mu}{\sigma} \leq 3\} - 1 = \fbox{0.9974} $
            \item   If you have data that you think follows a normal distribution, you can find the average to estimate $\mu$.
                    The calculations in (a)-(c) can similarly be used to estimate $\sigma$. Explain how you would go
                    about doing so.
                    \\Answer:\\
                    One can estimate the $\sigma$ values by assigning $Z$ values of 1, 2, or 3, and rearranging the 
                    equation to obtain $\sigma = |\frac{X - \mu}{Z}|$, with the absolute value taken to account for 
                    negative sigmas, and where $\mu$ is the median value of the data at the
                    50\% mark and $X$ is the value obtained at the points deviating from the median by $\pm$ 34.14\%, 47.72\%,
                    and 49.87\% for $\sigma$, $2\sigma$, and $3\sigma$ respectively. Each of these obtained sigmas may then be
                    averaged, given by 
                    $\sigma_{\text{avg}} = \frac{\sigma_{+1}+\sigma_{-1}+\sigma_{+2}+\sigma_{-2}+\sigma_{+3}+\sigma_{-3}}{6}$.
        \end{enumerate}
    \clearpage
    \item   Some time ago, a random sample of 747 obituaries published in Salt Lake City newspapers found that 46\% of 
            the decedents died in the 3-month period following their birthdays.
        \begin{enumerate}
            \item   What percent would you expect to die in that 3-month period? What assumption(s) are you making?\\
                Answer:\\
                Given the assumption that deaths are evenly distributed, it would follow that 25\% of deaths would occur
                in the 3-month period following a birthday, given that would equate to 25\% of the calendar year.
            \item   Estimate the probability that 46\% or more of the decedents would die in the 3-month period following 
                    their birthdays. \\
                Answer:\\
                Given the assumption $X \sim \text{Binomial}(747,.25)$, it may be approximated through a standard normal distribution, where
                \begin{align*}
                    46\% \times 747 &= 343.62 \\
                    P(X \geq 343.62) &= 1 - P(X < 343.62) \\
                    E[X] &= np = 186.75 \\
                    Var(X) &= np(1-p) = 186.75(.75) = 140.0625\\
                    \sigma &= \sqrt{Var(X)} \approx 11.8348
                \end{align*}
                Then, given that $ \frac{X - \mu}{\sigma} = Z \sim \text{normal}(0,1) $,\\ it follows that
                $$ \frac{343.62-186.75}{11.8348} = Z \approx 13.25$$
                Since the $Z$ value obtained is roughly 4 times larger than the largest $Z$ value on the table provided,
                and since the cdf approaches 1 as $Z\rightarrow\infty$,
                it can be concluded that
                $$ P(X \geq 343.62) = 1 - P(X < 343.62) \approx 1 - 1 = \fbox{0} $$
            \item   What conclusion might you draw from your calculation in (b)? How many deaths in the 3-month period 
                    would you accept as consistent with your expectations?
                    \\Answer:\\
                    Conclusions may include:
                    \begin{enumerate}
                        \item   data collection or calculation error.
                        \item   clustering of birthdays with some mass death event in the sampled decedent data.
                        \item   underlying causal factor(s) in the mortality rate in the 3-month period following a birthday.
                    \end{enumerate}
                    In the expected even distribution case, it would be expected to have roughly 187 $\pm$ 12 deaths occur.
        \end{enumerate}
    \clearpage
    \item   Suppose $Y_j$, where $1 \leq j \leq n$, are independent exponential random variables with parameters 
            $\lambda_j$. Find the cdf and pdf for the random variable $M = \text{min}(Y_1,Y_2,\ldots,Y_n)$. 
            What type of random variable is $M$?
            \\Answer:\\
            Given the hint to find $P(M>m)$ for a given $M\geq0$,
            \begin{align*}
                P(M>m) &= P(Y_1 > m) \times \cdots \times P(Y_n > m) \\
                P(Y_j > m) &= 1 - P(Y_j < m) = 1 - F_Y(m) = e^{-\lambda m} \\
                \Rightarrow P(M>m) &= 1 - \prod_{i=1}^{n} e^{-\lambda_{i} m} =1 - exp\left(-m\sum_{i=1}^{n}\lambda_{i}\right)
            \end{align*}
            Then, assigning $\sum \lambda_{i} = \lambda_{M}$ gives a cdf of,
            $$ \fbox{$1 - e^{-m\lambda_{M}}$} $$
            with the pdf given by
            $$F'_M(m) = f_M(m) = 1 - e^{-m\lambda_M} \frac{d}{dm} =\fbox{$ \lambda_{M} e^{-m\lambda_{M}}$} $$
            resulting in M being an exponential random variable, given its cdf and pdf match the general forms of each.
    \clearpage
    \item   Suppose that an appliance is constructed in such a way that it requires that $n$ independent electronic 
            components are all functioning. Assume that the lifespan of each of the components, $T_j$, is an exponential 
            random variable with parameter $\lambda_j$.
        \begin{enumerate}
            \item   Let $X$ be the random variable giving the lifespan of the appliance. Find the distribution of $X$. 
                \\Answer:\\
                As was obtained in question 3, lifespan $X$ would be an exponential random variable, given that the appliance 
                will break when one of its components break (or, in other words, the minimum $T_j$), and all components 
                $T_j$ are exponential random variables. 
            \item   Find $E[X]$. Then, find the \textbf{median} lifespan of the appliance, that is, the time $t$ at 
                    which half of the appliances are likely to have broken and half still working.
                    \\Answer:\\
                    Given X is an exponential r. v., it follows that 
                    $$ E[X] = \frac{1}{\lambda_{X}}, \lambda_{X} = \sum_{i=1}^{n} \lambda_{i} $$
                    with the median being calculated by
                    $$ F_X(t) = 1 - e^{-\lambda_X t} \Rightarrow \frac{1}{2} = 1 - e^{-\lambda_X t} $$
                    $$ 
                        \frac{1}{2} = e^{-\lambda_X t} \Rightarrow t = -\ln\left(\frac{1}{2}\right)/\lambda_X = 
                        -\ln\left(\frac{1}{2}\right)E[X]
                    $$
            \item   Consider the two numbers you computed in (b). What does this indicate in context about the 
                    lifespans of these devices? Which will the manufacturer of the appliance likely use in their advertising?
                    \\Answer:\\
                    Given that $-\ln\left(\frac{1}{2}\right) \approx 0.693$, the average appliance will last longer 
                    than 50\% of other such appliances. This would lead to the manufacturer likely advertising the 
                    average lifespan of the appliance.
        \end{enumerate}
    \clearpage
    \item   Let $X \sim exponential(\lambda)$, and let $Y = e^X$.
        \begin{enumerate}
            \item   Find the cdf of $Y$.
                \\Answer:
                \begin{align*}
                    F_Y(y) &= P(Y\leq y) = P(e^X \leq y) = P(X \leq \ln(y))\\
                           &= 1 - e^{-\lambda \ln(y)} = \fbox{$1 - y^{-\lambda}$}
                \end{align*}
            \item   Use (a) to find the pdf of $Y$. Check that your answer really is a density by showing 
                    that it integrates to 1.
                \\Answer:
                $$ F'_Y(y) = f_Y(y) = 1 - y^{-\lambda} \frac{d}{dy} = \fbox{$\lambda y^{-\lambda}$}$$
                Then, since the smallest possible value $X$ can take is 0, then the smallest $Y$ is $e^0 = 1$, 
                giving limits of integration of $[1,\infty)$
                $$ 
                    \int^{\infty}_{1} \lambda y^{-\lambda -1} dy = 
                    {-y^{-\lambda}}\bigg|^{\infty}_{1} = 0 - \left(-\frac{1}{1^\lambda}\right) = \fbox{1}
                $$
            \item   Find the pdf of Y again, this time using the "shortcut" formula (the Theorem from Part 3 of the Ch 5 Notes), 
                    being sure to explain why it applies.
        \end{enumerate}
\end{enumerate}

\end{document}
